\section{Mathematical Foundations}
\label{app:mathematical-foundations}
This appendix supplies citation anchors for the analytic notation used in the
main chapters and in Appendix~\ref{app:mathematical-notation}. It separates domain
and boundary language, classical function spaces, field notation, differential
operators, and norm conventions so later sections can cite only the convention
they need.

\begin{convention}[Standing regularity convention]
\label{conv:foundation-standing-regularity}
Unless a section states a stronger or weaker hypothesis, classical identities are
used only under the regularity needed for their terms to be defined. In
particular, fields have the displayed classical derivatives in the indicated
domains, boundary traces are defined when boundary terms are written, and
integration-by-parts or adjoint identities are invoked only when the required
boundary terms are meaningful. When a weak formulation is discussed, the relevant
Sobolev space, trace theorem, and boundary regularity assumptions are stated
explicitly; Sobolev traces are not inferred from this classical convention. At
minimum, boundaries are assumed piecewise $C^1$ when classical boundary
integrals or integration-by-parts identities are used. Outward unit normals are
then interpreted on the regular boundary pieces, hence surface-a.e. When a
pointwise normal derivative is asserted on all of $\pOmg$, the boundary is
assumed at least $C^1$.
\end{convention}

\begin{definition}[Domain and boundary notation]
\label{def:foundation-domain-boundary-notation}
Let $\Omg \subset \R^n$ be a bounded open connected set, with $n\in\N$, and
denote its closure by $\OmgBar$. In the hemodynamics settings considered in
this report, only $n\in[3]$ is used. The standing regularity convention in
Convention~\ref{conv:foundation-standing-regularity} applies whenever traces,
outward unit normals, or integration by parts are used. For open $\Omg$,
\begin{flalign*}
    \quad \OmgBar=\Omg\cup\pOmg, \qquad \pOmg=\OmgBar\setminus\Omg. &&
\end{flalign*}
Since $\Omg$ is bounded, $\OmgBar$ is closed and bounded and is therefore
compact by the Heine--Borel theorem. The open set $\Omg$ itself need not be
compact. Since $\Omg$ is bounded, $\pOmg$ is also compact. When the standing
regularity convention supplies a piecewise $C^1$ boundary, the outward unit
normal $\nhat$ is interpreted on the regular boundary pieces, hence $\dSx$-a.e.
on $\pOmg$.
\end{definition}

\begin{definition}[Time intervals and compact space-time sets]
\label{def:foundation-time-and-compact-set-notation}
For a final time $T_{\mathrm{final}}>0$, this report uses
\begin{flalign*}
    \quad I_T \defined (0,T_{\mathrm{final}}),
    \qquad
    \ITBar \defined [0,T_{\mathrm{final}}]. &&
\end{flalign*}
The open interval $I_T$ is used for interior-time differential identities and
space-time PDE regions, while $\ITBar$ is used when initial-time data,
endpoint data, sampled time grids, or flow maps are included. When a local final
time such as $T_\ast$ is stated, the same convention applies with
$T_{\mathrm{final}}$ replaced by that local final time.

For one-dimensional output fields, a compact space-time rectangle is written as
\begin{flalign*}
    \quad K \defined [0,\ell]\times[0,T_{\mathrm{final}}] &&
\end{flalign*}
or by an explicitly stated local variant. Spaces such as $C(K)$ and
$C^1([0,\ell])$ mean continuous functions on the named compact set, with the
supremum norm when a norm is written. Spaces such as $L^p(K)$ use the product
Lebesgue measure and the same almost-everywhere convention as
Definition~\ref{def:foundation-norm-conventions}.
\end{definition}

\begin{definition}[Field and multi-index notation]
\label{def:foundation-field-and-derivative-notation}
For $\Omg\subset\R^n$, we use the following field notation:
\begin{flalign*}
    \quad & \phi:\Omg\to\R \;\text{ is a \emph{scalar field}}, & \\
    \quad & \vf:\Omg\to\R^{n} \;\text{ is a \emph{vector field}}, &&\\
    \quad & \vT:\Omg\to\R^{n\times n} \;\text{ is a \emph{(second-order) tensor field}},&&
\end{flalign*}
and we write $\phi(\vx)$, $\vf(\vx)$, and $\vTx$\footnote{
A second-order tensor value $\vTx\in\R^{n\times n}$ can be identified with
the matrix of a linear map $\R^n\to\R^n$ acting on vectors by matrix-vector
multiplication.}
for the values at a point $\vx \in \Omg$. When referring to a physical quantity,
its measurement units are indicated as $\units{\cdot}$, whereas dimensionless
quantities are indicated as $\units{-}$. The $k$-th partial derivative of
$\phi$ with respect to coordinate $x_i$ is
\begin{flalign*}
    \quad & \pd_{x_i}^k \phi \equivto \frac{\pd^k \phi}{\pd x_i^k}. &&
\end{flalign*}
For a multi-index $\alpha=(\alpha_1,\ldots,\alpha_n)\in\N_0^n$, set
$|\alpha|=\alpha_1+\cdots+\alpha_n$. When the derivative exists classically,
the corresponding derivative of a scalar field $\phi$ is
\begin{flalign*}
    \quad & D^\alpha \phi \equivto \frac{\pd^{|\alpha|} \phi}{\pd x_1^{\alpha_1} \, \pd x_2^{\alpha_2} \, \ldots \, \pd x_n^{\alpha_n}}. &&
\end{flalign*}
For $|\alpha| = 1$, $D^\alpha \phi$ is a first-order partial derivative of
$\phi$, e.g. $\pd_{x_i} \phi$
where $\alpha_i = 1$ and $\alpha_j = 0$ for all $j\in[n]\setminus\{i\}$.
\end{definition}

\begin{definition}[Classical function-space notation]
\label{def:foundation-classical-function-space-notation}
Given a domain $\Omg$ and a scalar field $\mathbb F\in\{\R,\C\}$, a scalar
function space on $\Omg$ is an $\mathbb F$-vector subspace of the maps
$f:\Omg\to \mathbb F$. Unless a complex-valued space is
explicitly specified, this report uses real-valued spaces. Thus $C^0(\Omg)$
denotes the real-valued continuous functions on the open domain $\Omg$.
For $k\in\N_0$, $C^k(\Omg)$ denotes functions $\phi:\Omg\to\R$ such that
$D^\alpha\phi$ exists and is continuous in $\Omg$ for every $|\alpha|\le k$.
By contrast, $C^k(\OmgBar)$ denotes functions in $C^k(\Omg)$ whose derivatives
through order $k$ extend continuously to the closure $\OmgBar$. The space
$C^\infty(\Omg)$ is the intersection of all $C^k(\Omg)$ for $k\in\N_0$, i.e.
the infinitely differentiable functions on $\Omg$.
The notation $C^k(\Omg;\R^m)$ denotes the maps $\vf:\Omg \to \R^m$ with
components in $C^k(\Omg)$, i.e., $\vf = (f_i)_{i\in[m]}$ with
$f_i \in C^k(\Omg)$ for all $i \in[m]$.
\end{definition}

\begin{definition}[Differential-operator conventions]
\label{def:foundation-differential-operator-conventions}
For real-valued $\phi \in C^1(\Omg)$, the partial derivative with respect to coordinate
$x_i$ is
\begin{flalign*}
    \quad \pd_{x_i} \phi = \frac{\pd \phi}{\pd x_i}. &&
\end{flalign*}
Named-coordinate partial derivatives use the same compact convention; for
example,
\begin{flalign*}
    \quad
    \pd_t \phi = \frac{\pd \phi}{\pd t},
    \qquad
    \pd_z f = \frac{\pd f}{\pd z},
    \qquad
    \pd_A p = \frac{\pd p}{\pd A}. &&
\end{flalign*}
Cylindrical derivatives such as $\pd_r u$ and $\pd_\theta u$ are read
analogously. Boundary notation retains $\partial$, as in $\partial V(t)$ and
$\pOmg$.
For $\vx \in \Omg$ and a direction $\vv \in \R^n$, the directional derivative is
\begin{flalign*}
    \quad D_{\vv} \phi(\vx) \defined \inp{\nabla \phi(\vx)}{\vv}_{\R^n}. &&
\end{flalign*}
If $\phi\in C^1(\OmgBar)$, $\vx\in\pOmg$, and $\nhat$ is the outward unit
normal at $\vx$, then the boundary normal derivative is
\begin{flalign*}
    \quad \pd_{\nhat} \phi(\vx) \defined \inp{\nabla \phi(\vx)}{\nhat}_{\R^n}. &&
\end{flalign*}
For Lebesgue measure $\dx$ on $\R^n$ the integral of $\phi$ in $\Omg$ is
\begin{flalign*}
    \quad \int_{\Omg} \phi(\vx) \dx. &&
\end{flalign*}
If $\pOmg$ is of class $C^1$ and $\phi\in C^0(\OmgBar)$, then the classical boundary trace
$\phi|_{\pOmg}$ is continuous on $\pOmg$. Since
$\Omg$ is bounded and $\pOmg$ is a compact $C^1$ hypersurface,
this trace belongs to $L^1(\pOmg; \dSx)$. The corresponding boundary integral
is written
\begin{flalign*}
    \quad \int_{\pOmg} \phi(\vx)\,\dSx. &&
\end{flalign*}
We use the column-gradient convention for scalar fields. For
$\phi\in C^1(\Omg)$, the gradient is
\begin{flalign*}
    \quad \nabla \phi \defined \left(\pd_{x_i} \phi\right)_{i\in[n]} \in \R^n. &&
\end{flalign*}
For $\phi\in C^2(\Omg)$, the Laplacian is
\begin{flalign*}
    \quad \del \phi \defined \nabla\cdot\nabla\phi \equals \sum_{i\in[n]} \pd_{x_i}^2 \phi. &&
\end{flalign*}

For vector-valued $\vf \in C^1(\Omg; \R^n)$ with components $\vf = (f_i)_{i\in[n]}$,
the vector gradient is the Jacobian with entries $(\nabla \vf)_{ij} \defined \pd_{x_j} f_i$.
The divergence is the trace of this Jacobian:
\begin{flalign*}
    \quad \Div \, \vf \defined \nabla \cdot \vf \equals \sum_{i\in[n]} \pd_{x_i} f_i
    \equals \operatorname{tr}(\nabla \vf). &&
\end{flalign*}
For $\vf \in C^2(\Omg; \R^n)$, the Laplacian is defined componentwise. The domain integral is
also componentwise whenever the components are integrable:
\begin{flalign*}
    \quad &  \del \, \vf \defined \left( \del f_i \right)_{i\in[n]} &&\\
    \quad &  \int_\Omg \vf(\vx)\,\dx \defined \left(\int_\Omg f_i(\vx)\,\dx\right)_{i\in[n]}. &&
\end{flalign*}

For tensor-valued $\vT \in C^1(\Omg; \R^{n\times n})$ with entries $\vT = (T_{ij})_{i,j\in[n]}$,
we use the row-divergence convention $(\Divbf\vT)_i=\sum_{j\in[n]} \pd_{x_j} T_{ij}$.
Equivalently, if $\vt_i=(T_{ij})_{j\in[n]}$ denotes the $i$-th row of $\vT$, then
\begin{flalign*}
    \quad \Divbf \, \vT &\defined \left( \Div \, \vt_{i} \right)_{i\in[n]} &&\\
        & \equals \left( \nabla \cdot \vt_{i} \right)_{i\in[n]} &&\\
        & \equals \left( \sum_{j\in[n]} \pd_{x_j} T_{ij} \right)_{i\in[n]}. &&
\end{flalign*}
For $\vT \in C^2(\Omg; \R^{n\times n})$, the Laplacian is defined componentwise.
The domain integral is also defined componentwise whenever the entries are integrable:
\begin{flalign*}
    \quad & \del \, \vT \defined \left( \del T_{ij} \right)_{i,j\in[n]} &&\\
    \quad & \int_\Omg \vT(\vx)\,\dx \defined \left( \int_\Omg T_{ij}(\vx)\,\dx \right)_{i,j\in[n]}. &&
\end{flalign*}
\end{definition}

\begin{definition}[$C^k$ and $L^p$ norm conventions]
\label{def:foundation-norm-conventions}
Let $\Omg\subset\R^n$ be bounded and open. We write
$C^k(\OmgBar)$ for the functions $\phi\in C^k(\Omg)$ such that
each derivative $D^\alpha\phi$, $|\alpha|\le k$, extends continuously to
$\OmgBar$. The norm is
\begin{flalign*}
    \quad \norm{\phi}_{C^k(\OmgBar)} \defined
    \sum_{|\alpha|\le k} \sup_{\vx\in\OmgBar} |D^\alpha\phi(\vx)|. &&
\end{flalign*}
With this norm, $C^k(\OmgBar)$ is a Banach space: every Cauchy
sequence in $C^k(\OmgBar)$ converges, with respect to $\norm{\cdot}_{C^k(\OmgBar)}$,
to a limit in $C^k(\OmgBar)$. In particular, for $k = 0$, we have
\begin{flalign*}
    \quad \norm{\phi}_{C^0(\OmgBar)} &\defined \sup_{\vx \in \OmgBar} \left| \phi(\vx) \right|
    = \max_{\vx \in \OmgBar} \left| \phi(\vx) \right|. &&
\end{flalign*}
For continuous functions on compact sets, the max norm agrees with the usual
supremum norm.

For $1\le p<\infty$, $L^p(\Omg)$ denotes the space of equivalence classes of
real-valued measurable functions on $\Omg$, identified up to equality almost
everywhere, with finite norm
\begin{flalign*}
    \quad \norm{\phi}_{L^p(\Omg)} \defined \left(\int_\Omg |\phi(\vx)|^p \dx\right)^{1/p}. &&
\end{flalign*}
For $p=\infty$,
\begin{flalign*}
    \quad \norm{\phi}_{L^\infty(\Omg)} \defined \operatorname*{ess\,sup}_{\vx\in\Omg}|\phi(\vx)|. &&
\end{flalign*}
With this almost-everywhere identification, $L^p(\Omg)$ is a Banach space for
$1\le p\le\infty$.

The notation $L^p(\Omg;\R^m)$ denotes equivalence classes of measurable maps
$\vf:\Omg\to\R^m$, identified up to equality a.e., whose components belong to
$L^p(\Omg)$, i.e., $\vf = (f_i)_{i\in[m]}$ with $f_i \in L^p(\Omg)$ for all $i \in[m]$.
When a norm on $L^p(\Omg;\R^m)$ is needed, we use the Euclidean norm in the
range. For $1\le p<\infty$,
\begin{flalign*}
    \quad
    \norm{\vf}_{L^p(\Omg;\R^m)}
    \defined
    \left(\int_\Omg \norm{\vf(\vx)}_{\R^m}^p\dx\right)^{1/p}.
    &&
\end{flalign*}
For $p=\infty$, the corresponding convention is
\begin{flalign*}
    \quad
    \norm{\vf}_{L^\infty(\Omg;\R^m)}
    \defined
    \operatorname*{ess\,sup}_{\vx\in\Omg}\norm{\vf(\vx)}_{\R^m}.
    &&
\end{flalign*}
For matrix or second-order tensor values, the pointwise range norm is the
Frobenius norm
\begin{flalign*}
    \quad \norm{\vT}_{\mathrm F}
    \defined \left(\vT:\vT\right)^{1/2}
    = \left(\sum_{i,j\in[n]}T_{ij}^{2}\right)^{1/2}. &&
\end{flalign*}
Thus tensor-valued $L^p$ norms are formed by replacing $|\phi(\vx)|$ in the
scalar $L^p$ formula with $\norm{\vT(\vx)}_{\mathrm F}$.
\end{definition}
